% GENERATED FILE. Edit canonical lessons, then run npm run generate:books.
% canonical-source-hash: fnv1a64:3fc8777b00275f39
% canonical-lessons: JA-C378-man-in, JA-C378-konzatsu, JA-C378-gyuugyuu, JA-C378-tsumekomu, JA-C378-tachiagaru

\chapter{A Crowded Carriage}
\label{ch:ja-a-crowded-carriage}
\hlchaptermodality{378}

\begin{chapteropening}
\textbf{By the end of this chapter:} \emph{I can describe a crowded train.}

Complete the last lesson of chapter 378: I can describe a crowded train.
\end{chapteropening}

\section[\texorpdfstring{\ja{まんいん} (man'in)}{まんいん (man'in)}]{\ja{まんいん} (man'in) — full (of people)}

\label{lesson:JA-C378-man-in}



\begin{quote}
Before the new one: say the Japanese for a direction (of travel), then the Japanese for to move away.
\end{quote}

\subsection*{You'll want to know: \ja{まんいん}}
\textbf{\ja{まんいん}} — \emph{man'in} — \textquotedblleft{}full (of people)\textquotedblright{}.

\begin{quote}
\textbf{\ja{あさの} \ja{でんしゃは} \ja{いつも} \ja{まんいんです。}} — \emph{asa no densha wa itsumo man'in desu} — \textquotedblleft{}The morning trains are always full.\textquotedblright{}
\end{quote}

\begin{cousinweb}
\textbf{\ja{まんいん}} (\emph{man'in}) is Sino-Japanese (kango): two Chinese characters read with their on'yomi, readings borrowed from Chinese. They are \emph{man}, \textquotedblleft{}full\textquotedblright{} (Middle Chinese \emph{manX}) and \emph{in}, \textquotedblleft{}member\textquotedblright{} (\emph{hjwen}) (KANJIDIC2; Middle Chinese in Baxter's transcription).
\end{cousinweb}

\begin{grammarlens}[title={What you've built}]
One more word for getting around.
\end{grammarlens}

\subsection*{Your turn}
\begin{itemize}
\raggedright
  \item \emph{Say it:} \emph{man'in}
  \item \emph{Say it:} \emph{man'in}, once more
  \item \emph{Say it:} \emph{asa no densha wa itsumo man'in desu} — The morning trains are always full.
  \item \emph{Recall:} say the Japanese for a direction (of travel), then the Japanese for to move away, then say \emph{man'in} again
\end{itemize}

\begin{tcolorbox}[breakable,colback=teal!4,colframe=teal!35!black,title={Before you move on}]
What does \ja{まんいん} mean? (\textquotedblleft{}Full (of people)\textquotedblright{}.) Say it once more.
\end{tcolorbox}
\section[\texorpdfstring{\ja{こんざつ} (konzatsu)}{こんざつ (konzatsu)}]{\ja{こんざつ} (konzatsu) — congestion; crowding}

\label{lesson:JA-C378-konzatsu}



\begin{quote}
Before the new one: say the Japanese for to move away, then the Japanese for full (of people).
\end{quote}

\subsection*{You'll want to know: \ja{こんざつ}}
\textbf{\ja{こんざつ}} — \emph{konzatsu} — \textquotedblleft{}congestion; crowding\textquotedblright{}.

\begin{quote}
\textbf{\ja{しゅうまつの} \ja{えきは} \ja{こんざつしていました。}} — \emph{shūmatsu no eki wa konzatsu shite imashita} — \textquotedblleft{}The station was crowded at the weekend.\textquotedblright{}
\end{quote}

\begin{cousinweb}
\textbf{\ja{こんざつ}} (\emph{konzatsu}) is Sino-Japanese (kango): two Chinese characters read with their on'yomi, readings borrowed from Chinese. They are \emph{kon}, \textquotedblleft{}mix\textquotedblright{} (Middle Chinese \emph{hwonX}) and \emph{zatsu}, \textquotedblleft{}miscellaneous\textquotedblright{} (\emph{dzop}) (KANJIDIC2; Middle Chinese in Baxter's transcription).
\end{cousinweb}

\begin{grammarlens}[title={What you've built}]
One more word for getting around.
\end{grammarlens}

\subsection*{Your turn}
\begin{itemize}
\raggedright
  \item \emph{Say it:} \emph{konzatsu}
  \item \emph{Say it:} \emph{konzatsu}, once more
  \item \emph{Say it:} \emph{shūmatsu no eki wa konzatsu shite imashita} — The station was crowded at the weekend.
  \item \emph{Recall:} say the Japanese for to move away, then the Japanese for full (of people), then say \emph{konzatsu} again
\end{itemize}

\begin{tcolorbox}[breakable,colback=teal!4,colframe=teal!35!black,title={Before you move on}]
What does \ja{こんざつ} mean? (\textquotedblleft{}Congestion; crowding\textquotedblright{}.) Say it once more.
\end{tcolorbox}
\section[\texorpdfstring{\ja{ぎゅうぎゅう} (gyūgyū)}{ぎゅうぎゅう (gyūgyū)}]{\ja{ぎゅうぎゅう} (gyūgyū) — packed tight; squeezed in}

\label{lesson:JA-C378-gyuugyuu}



\begin{quote}
Before the new one: say the Japanese for full (of people), then the Japanese for congestion.
\end{quote}

\subsection*{You'll want to know: \ja{ぎゅうぎゅう}}
\textbf{\ja{ぎゅうぎゅう}} — \emph{gyūgyū} — \textquotedblleft{}packed tight; squeezed in\textquotedblright{}.

\begin{quote}
\textbf{\ja{でんしゃは} \ja{ぎゅうぎゅうでした。}} — \emph{densha wa gyūgyū deshita} — \textquotedblleft{}The train was packed tight.\textquotedblright{}
\end{quote}

\begin{cousinweb}
\textbf{\ja{ぎゅうぎゅう}} (\emph{gyūgyū}) is a native mimetic word that pictures people or things pressed tightly together: an imitative Japanese form, inherited rather than borrowed, whose older origin is not settled (JMdict).
\end{cousinweb}

\begin{grammarlens}[title={What you've built}]
One more word for getting around.
\end{grammarlens}

\subsection*{Your turn}
\begin{itemize}
\raggedright
  \item \emph{Say it:} \emph{gyūgyū}
  \item \emph{Say it:} \emph{gyūgyū}, once more
  \item \emph{Say it:} \emph{densha wa gyūgyū deshita} — The train was packed tight.
  \item \emph{Recall:} say the Japanese for full (of people), then the Japanese for congestion, then say \emph{gyūgyū} again
\end{itemize}

\begin{tcolorbox}[breakable,colback=teal!4,colframe=teal!35!black,title={Before you move on}]
What does \ja{ぎゅうぎゅう} mean? (\textquotedblleft{}Packed tight; squeezed in\textquotedblright{}.) Say it once more.
\end{tcolorbox}
\section[\texorpdfstring{\ja{つめこむ} (tsumekomu)}{つめこむ (tsumekomu)}]{\ja{つめこむ} (tsumekomu) — to cram in; to stuff}

\label{lesson:JA-C378-tsumekomu}



\begin{quote}
Before the new one: say the Japanese for congestion, then the Japanese for packed tight.
\end{quote}

\subsection*{You'll want to know: \ja{つめこむ}}
\textbf{\ja{つめこむ}} — \emph{tsumekomu} — \textquotedblleft{}to cram in; to stuff\textquotedblright{}.

\begin{quote}
\textbf{\ja{かばんに} \ja{にもつを} \ja{つめこみました。}} — \emph{kaban ni nimotsu o tsumekomimashita} — \textquotedblleft{}I crammed my things into the bag.\textquotedblright{}
\end{quote}

Polite \textbf{\ja{つめこみます}}, \emph{tsumekomimasu}; a godan verb. Past (polite): \textbf{\ja{あねが} \ja{かばんに} \ja{ふくを} \ja{つめこみました。}}, \emph{ane ga kaban ni fuku o tsumekomimashita}.

\begin{cousinweb}
\textbf{\ja{つめこむ}} (\emph{tsumekomu}) is a native Japanese word (wago), inherited rather than borrowed from Chinese, built from \emph{tsume}, \textquotedblleft{}packing\textquotedblright{} (from \emph{tsumeru}) and \emph{komu}, \textquotedblleft{}to go in\textquotedblright{}; its kanji were attached later as kun'yomi (JMdict).
\end{cousinweb}

\begin{grammarlens}[title={What you've built}]
One more word for getting around.
\end{grammarlens}

\subsection*{Your turn}
\begin{itemize}
\raggedright
  \item \emph{Say it:} \emph{tsumekomu}
  \item \emph{Say it:} \emph{tsumekomu}, once more
  \item \emph{Say it:} \emph{kaban ni nimotsu o tsumekomimashita} — I crammed my things into the bag.
  \item \emph{Recall:} say the Japanese for congestion, then the Japanese for packed tight, then say \emph{tsumekomu} again
\end{itemize}

\begin{tcolorbox}[breakable,colback=teal!4,colframe=teal!35!black,title={Before you move on}]
What does \ja{つめこむ} mean? (\textquotedblleft{}To cram in; to stuff\textquotedblright{}.) Say it once more.
\end{tcolorbox}
\section[\texorpdfstring{\ja{たちあがる} (tachiagaru)}{たちあがる (tachiagaru)}]{\ja{たちあがる} (tachiagaru) — to stand up}

\label{lesson:JA-C378-tachiagaru}



\begin{quote}
Before the new one: say the Japanese for packed tight, then the Japanese for to cram in.
\end{quote}

\subsection*{You'll want to know: \ja{たちあがる}}
\textbf{\ja{たちあがる}} — \emph{tachiagaru} — \textquotedblleft{}to stand up\textquotedblright{}.

\begin{quote}
\textbf{\ja{おじいさんが} \ja{きたので、たちあがりました。}} — \emph{ojiisan ga kita node, tachiagarimashita} — \textquotedblleft{}An old man came, so I stood up.\textquotedblright{}
\end{quote}

Polite \textbf{\ja{たちあがります}}, \emph{tachiagarimasu}; a godan verb. Past (polite): \textbf{\ja{こどもが} \ja{たちあがりました。}}, \emph{kodomo ga tachiagarimashita}.

\begin{cousinweb}
\textbf{\ja{たちあがる}} (\emph{tachiagaru}) is a native Japanese word (wago), inherited rather than borrowed from Chinese, built from \emph{tachi}, \textquotedblleft{}standing\textquotedblright{} (from \emph{tatsu}) and \emph{agaru}, \textquotedblleft{}to rise\textquotedblright{}; its kanji were attached later as kun'yomi (JMdict).
\end{cousinweb}

\begin{grammarlens}[title={What you've built}]
One more word for getting around.
\end{grammarlens}

\subsection*{Your turn}
\begin{itemize}
\raggedright
  \item \emph{Say it:} \emph{tachiagaru}
  \item \emph{Say it:} \emph{tachiagaru}, once more
  \item \emph{Say it:} \emph{ojiisan ga kita node, tachiagarimashita} — An old man came, so I stood up.
  \item \emph{Recall:} say the Japanese for packed tight, then the Japanese for to cram in, then say \emph{tachiagaru} again
\end{itemize}

\begin{tcolorbox}[breakable,colback=teal!4,colframe=teal!35!black,title={Before you move on}]
What does \ja{たちあがる} mean? (\textquotedblleft{}To stand up\textquotedblright{}.) Say it once more.
\end{tcolorbox}
